\clearpage
\SourceLeaf{357}\pdfbookmark[1]{Nachlass und Veröffentlichung}{parabolisch-veroeffentlichung}
\begin{center}
{\large NACHLASS UND VERÖFFENTLICHUNG.\par}
\bigskip
\pdfbookmark[2]{II. Zur Berechnung der wahren Anomalie in der parabolischen Bewegung}{parabolisch-anomalie}
\textbf{[II. Zur Berechnung der wahren Anomalie in der parabolischen Bewegung.]}\par
\bigskip
\pdfbookmark[3]{1. Neue Tafel für die parabolische Bewegung}{tabula-nova}
\textbf{[1. Neue Tafel für die parabolische Bewegung.]}\par
\medskip
\textbf{Erläuterung der neuen Tafel für die parabolische Bewegung.}
\end{center}
\medskip
\pdfbookmark[4]{1.}{tabula-nova-001}
\begin{center}1.\end{center}
\noindent\hspace*{\parindent}Wie in Art. 18–20 der Theoria Motus bezeichnen wir mit \(q\) den Abstand von der Sonne im Perihel, mit \(v\) die wahre Anomalie, mit \(t\) die seit dem Durchgang durch das Perihel verstrichene Zeit und schließlich mit [\(k\)] die Konstante \(3548''{,}188\) (Art. [6]). Da bei der parabolischen Bewegung eines Kometen, dessen Masse wir vernachlässigen, die in Bogensekunden ausgedrückte Winkelgeschwindigkeit im Allgemeinen
\[
\frac{\mathrm{d}v}{\mathrm{d}t}=\frac{k\BareRoot{2}}{q^{\frac32}}\cos\tfrac12v^4,
\]
beträgt, setzen wir der Kürze halber
\[
\frac{k\BareRoot{2}}{q^{\frac32}}=A,
\]
sodass \(A\) die Winkelgeschwindigkeit im Perihel angibt, die man zweckmäßig den parabolischen Elementen eines jeden Kometen hinzufügt. Es gilt
\[
\log k\BareRoot{2}=3{,}700\ 5216.
\]
\par\medskip
\pdfbookmark[4]{2.}{tabula-nova-002}
\begin{center}2.\end{center}
\noindent\hspace*{\parindent}Der Zusammenhang zwischen der wahren Anomalie und der Zeit wird durch die Formel
\SourceLeaf{358}
\[
\operatorname{tang}\tfrac12v+\tfrac13\operatorname{tang}\tfrac12v^3
=\frac{At}{2\mathpunct{.}206265''}=\frac{\pi}{1296000''}At
\]
ausgedrückt, wobei \(\pi\) den halben Umfang eines Kreises mit Einheitsradius bezeichnet. Die Zeit ergibt sich somit unmittelbar aus der wahren Anomalie. Für die umgekehrte Rechnung liefert die \textsc{Cardano}sche Regel unter Einführung zweier Hilfswinkel \(P\), \(Q\) die folgenden Formeln:
\[
\operatorname{tang}P=\frac{4\mathpunct{.}206265''}{3At}=\frac{864000''}{\pi At}
\]
\[
\operatorname{tang}Q=\BareRoot[3]{\operatorname{tang}\tfrac12P}
\]
\[
\operatorname{tang}\tfrac12v=2\operatorname{cotang}2Q.
\]
\par\medskip
\pdfbookmark[4]{3.}{tabula-nova-003}
\begin{center}3.\end{center}
\noindent\hspace*{\parindent}Da aber die letztere Aufgabe sehr häufig vorkommt, wäre es mühsam, stets auf diese Formeln oder auf indirekte Verfahren zurückzugreifen. Daher verwenden die Astronomen zu diesem Zweck gewöhnlich eine Hilfstafel. Die gebräuchliche Tafel liefert zum Argument
\[
\frac{t}{q^{\frac32}}=\frac{At}{k\BareRoot{2}}
\]
die wahre Anomalie. Der berühmte \textsc{Burckhardt} wählte den Logarithmus dieser Größe\footnote{Connaiss. des tems 1816.}. Das Argument der \textsc{Barker}schen Tafel ist die wahre Anomalie, neben der der zugehörige Wert der Größe
\[
\frac{k\pi\BareRoot{2}\mathpunct{.}t}{17280''\mathpunct{.}q^{\frac32}}=\frac{\pi At}{17280''}
\]
angegeben ist. In der diesem Werk beigefügten neuen Tafel (die bereits erstellt war, als die Tafel des berühmten \textsc{Burckhardt} veröffentlicht wurde) haben wir den Logarithmus der Größe \(At\) als Argument gewählt und ihn in Schritten von einem Tausendstel der Einheit fortschreiten lassen. Da bei der parabolischen Bewegung der Wert von \(At\) von \(0\) bis ins Unendliche wächst, hätte die Tafel auf diese Weise offenbar weder am Anfang noch am Ende vollständig werden können. Deshalb haben wir am Anfang der Tafel die Größe \(At\) selbst als Argument genommen und sind in Hunderterschritten bis zum Wert \(40000\) fortgeschritten, dem die wahre Anomalie \(11^0\ 2'\ 32''\) entspricht. Dann laufen die Logarithmen von \(4{,}6\) bis \(8{,}2\). Diesen Grenzen entsprechen die Anomalien \(10^0\ 59'\ 26''\) und \(169^0\ 0'\ 15''\).
\par
\SourceLeaf{359}
Größere Anomalien werden nur sehr selten benötigt. Damit die Tafel dennoch an beiden Enden vollständig ist, haben wir einen dritten Teil hinzugefügt, in dem wir als Argument die Größe
\[
\BareRoot[3]{\frac{B}{t}}\mathpunct{\raisebox{.45ex}{,}}
\]
gewählt haben, wobei
\[
B=\frac{1296000^4}{3\pi^4A}=\frac{1296000^4}{3\pi^4k\BareRoot{2}}q^{\frac32}.
\]
\par
\Needspace{8\baselineskip}
Dieses Argument erstreckt sich absteigend vom Wert \(40000\) bis \(0\). Diesen Werten entsprechen die Anomalien \(168^0\ 49'\ 9''\) und \(180^0\). Zur Bestimmung des Logarithmus der Konstanten \(B\) geben wir die folgenden Logarithmen an
\[
\log\frac{1296000^4}{3\pi^4}=[22{,}461\ 8207]
\]
\[
\log\frac{1296000^4}{3\pi^4k\BareRoot{2}}=[18{,}761\ 2991.]
\]
\par
Der Anwendungsbereich des ersten Teils der Tafel reicht somit
\[
\begin{array}{rll}
\text{von}&t=0&\text{bis }[t=]\ 7{,}971\,q^{\frac32}\\
\text{der des zweiten Teils von}&t=7{,}934\,q^{\frac32}&\text{bis }t=31585\,q^{\frac32}\\
\text{der des dritten Teils von}&t=30061\,q^{\frac32}&\text{bis ins Unendliche.}
\end{array}
\]
\par
Im Übrigen bemerken wir, dass die ganze Tafel für die Anomalie ursprünglich auf Hunderttausendstel einer Bogensekunde berechnet wurde, sodass man sich stets auf die zweite Stelle verlassen konnte.
\par\medskip
\pdfbookmark[4]{4.}{tabula-nova-004}
\begin{center}4.\end{center}
\noindent\hspace*{\parindent}Es lohnt sich, den Anfang und das Ende unserer Tafel getrennt zu betrachten und die Anomalie in eine nach Potenzen des Arguments fortschreitende Reihe zu entwickeln.
\par
Wir setzen
\[
\operatorname{tang}\tfrac12v=\theta
\]
\[
\frac{At}{206265''}=\alpha
\]
und erhalten dann
\[
\theta=\tfrac12\alpha-\tfrac13\theta^3.
\]
Daraus ergibt sich mithilfe des \textsc{Lagrange}schen Satzes im Allgemeinen
\SourceLeaf{360}
\[
\begin{aligned}
\theta^n={}&\left(\frac{\alpha}{2}\right)^n-\frac{n}{3}\left(\frac{\alpha}{2}\right)^{n+2}
+\frac{n\mathpunct{.}n+5}{3\mathpunct{.}6}\left(\frac{\alpha}{2}\right)^{n+4}\\
&-\frac{n\mathpunct{.}+n+7\mathpunct{.}n+8}{3\mathpunct{.}6\mathpunct{.}9}\left(\frac{\alpha}{2}\right)^{n+6}\\
&+\frac{n\mathpunct{.}n+9\mathpunct{.}n+10\mathpunct{.}n+11}{3\mathpunct{.}6\mathpunct{.}9\mathpunct{.}12}\left(\frac{\alpha}{2}\right)^{n+8}-\text{usw.}
\end{aligned}
\]
Setzen wir diesen Ausdruck für die einzelnen Potenzen \(\theta^n\) in die Reihe
\[
v=2\theta-\tfrac23\theta^3+\tfrac25\theta^5-\tfrac27\theta^7+\text{usw.}
\]
ein, so erhalten wir
\[
\begin{aligned}
v={}&\alpha-\tfrac13\cdot\tfrac14\left(1+\tfrac33\right)\alpha^3\\
&+\tfrac15\cdot\tfrac1{16}\left(1+\tfrac53+\tfrac{5\mathpunct{.}6}{3\mathpunct{.}6}\right)\alpha^5\\
&-\tfrac17\cdot\tfrac1{64}\left(1+\tfrac73+\tfrac{7\mathpunct{.}8}{3\mathpunct{.}6}+\tfrac{7\mathpunct{.}8\mathpunct{.}9}{3\mathpunct{.}6\mathpunct{.}9}\right)\alpha^7\\
&+\tfrac19\cdot\tfrac1{256}\Bigl(1+\tfrac93+\tfrac{9\mathpunct{.}10}{3\mathpunct{.}6}+\tfrac{9\mathpunct{.}10\mathpunct{.}11}{3\mathpunct{.}6\mathpunct{.}9}\\
&\hspace{5em}+\tfrac{9\mathpunct{.}10\mathpunct{.}11\mathpunct{.}12}{3\mathpunct{.}6\mathpunct{.}9\mathpunct{.}12}\Bigr)\alpha^9-\text{usw.,}
\end{aligned}
\]
Das Bildungsgesetz dieser Reihe liegt auf der Hand. Bestimmt man die Zahlenwerte der Koeffizienten, so erhält man
\[
\begin{aligned}
v={}&\alpha-\tfrac16\alpha^3+\tfrac{13}{240}\alpha^5-\tfrac{43}{2016}\alpha^7\\
&+\tfrac{191}{20736}\alpha^9-\tfrac{479}{114048}\alpha^{11}+\text{usw.}
\end{aligned}
\]
\par
Offenbar ist \(v\) hier als ein in Teilen des Radius ausgedrückter Bogen aufzufassen. Wird er in Bogensekunden verlangt, so ist die Reihe mit \(206265''\) zu multiplizieren, sodass sich ergibt
\[
v=At-\tfrac16\frac{(At)^3}{206265^2}+\tfrac{13}{240}\frac{(At)^5}{206265^4}-\text{usw.}
\]
Kleine Anomalien weichen also nur wenig von den zugehörigen Argumenten ab.
\par
Im Übrigen liefert der oben angegebene Ausdruck für \(\theta^n\) auch eine elegante Reihe für den Radiusvektor \(r\). Es ergibt sich nämlich
\[
\begin{aligned}
r={}&q(1+\theta\theta)\\
={}&q\Bigl(1+\tfrac14\alpha\alpha-\tfrac13\cdot\frac{\alpha^4}{8}+\frac{7}{3\mathpunct{.}6}\frac{\alpha^6}{32}\\
&-\frac{9\mathpunct{.}10}{3\mathpunct{.}6\mathpunct{.}9}\frac{\alpha^8}{128}
+\frac{11\mathpunct{.}12\mathpunct{.}13}{3\mathpunct{.}6\mathpunct{.}9\mathpunct{.}12}\frac{\alpha^{10}}{512}\\
&-\frac{13\mathpunct{.}14\mathpunct{.}15\mathpunct{.}16}{3\mathpunct{.}6\mathpunct{.}9\mathpunct{.}12\mathpunct{.}15}\frac{\alpha^{12}}{2048}+\text{usw.}\Bigr).
\end{aligned}
\]
Es lässt sich zeigen, dass beide Reihen konvergieren, solange \(\alpha<\tfrac43\) oder, was dasselbe ist, \(v<61^0\ 35'\ 45''{,}4\).
\par
